\begin{description}
\item[a)] Total energy of the projectile is
\[
  E = \tfrac{1}{2}mv_\circ^2 - \frac{GMm}{R} = -\frac{GMm}{2R} < 0
\]
$E < 0$ means that orbit might be ellipse or circle. As $\theta > 0$, the
orbit is an ellipse. Total energy for an ellipse is
\[
  E = -\frac{GMm}{2a}
\]
Then
\[
  a = R
\]
\hfill(7 points)

\item[b)] In figure (1) we have
\begin{align*}
  OA + O'A &= 2a \\
  O'A &= a
\end{align*}
% FIGURE NEEDED: Figure (1) - circle of the planet (centre O, radius R = a) and the orbit ellipse through launch points A (top) and B (bottom); horizontal major axis L-O-C-M-O'-H, with OA = O'A = OB = a marked, v_0 vector at A parallel to the major axis (2009 TH solutions booklet, p. 56)
In $OAO'$ triangle it is obvious that
\[
  OC = CO'
\]
Then $C$ must be the center of the ellipse with the initial velocity vector
$v_\circ$ parallel to the ellipse major-axis $(LH)$.

In figure (2)
\[
  HM = CH - CM = a - (R - R\sin\theta) = R - R + R\sin\theta
     = R\sin\theta = \frac{R}{2}
\]
\hfill(15 points)

\item[c)] Range of the projectile is $\overset{\frown}{AB}$
\[
  \overset{\frown}{AB} = 2\left(\frac{\pi}{2} - \theta\right)R
  = (\pi - 2\theta)R = \frac{2\pi}{3}R
\]
\hfill(6 points)

\item[d)]
% FIGURE NEEDED: Figure (2) - same construction as Figure (1) with the launch angle theta at A between v_0 and the tangent, angle pi/2 - theta at O and theta at M marked (2009 TH solutions booklet, p. 57)
Start with ellipse equation in polar coordinates
\[
  r = \frac{a(1 - e^2)}{1 + e\cos\varphi}
\]
For point A
\[
  R = \frac{R(1 - e^2)}{1 - e\cos(\frac{\pi}{2} + \theta)}
\]
\[
  e = \sin\theta = \frac{1}{2}
\]
\hfill(5 points)

\item[e)] Using Kepler's second law
% FIGURE NEEDED: ellipse with the swept region AOBH shaded grey (triangle AOB plus the half-ellipse from A over H to B), semi-axes a and b marked at centre C (2009 TH solutions booklet, p. 58)
\begin{align*}
  \frac{\Delta S}{S_0} &= \frac{\Delta T}{T} \\
  \Delta S = S_{AOBH} &= S_{\Delta AOB} + \frac{S_0}{2} \\
  &= 2\times\frac{bae}{2} + \frac{\pi ab}{2}
   = bae + \frac{\pi ab}{2} \\
  \frac{\Delta S}{S} &= \frac{bae + \frac{\pi ab}{2}}{\pi ab}
   = \frac{e + \frac{\pi}{2}}{\pi} = \frac{0.5 + \frac{\pi}{2}}{\pi}
\end{align*}
Kepler's third law
\[
  T = \sqrt{\frac{4\pi^2 R^3}{GM}} = 84.5\ min
\]
\[
  \Delta T = T\times\frac{0.5 + \frac{\pi}{2}}{\pi} = 55.7\ min
\]
\hfill(12 points)
\end{description}
